\documentclass[
    a4paper,
    DIV=13,
    abstract=true,
    numbers=noenddot
]{scrartcl}

\usepackage{preamble}
\addbibresource{library.bib} % duplicate of preamble, tells VimTeX where entries live

\begin{document}
\title{Langevin diffusion  
}
\author{Liam Llamazares-Elias}
\date{\today}
\maketitle
\section{Introduction}
\section{Langevin diffsion}
The equation 
\begin{aligneq}\label{LD}
    \d X_t=- \nabla \log{\pi}(X_t)+\sqrt{2} \d W_t.
\end{aligneq}

\section{Convergence}
\begin{theorem}
   The Langevin diffusion \eqref{LD} converges to $\pi$. 
\end{theorem}
\begin{proof}
    Direct substitution in the FP equation shows that $\pi$ is a statinoary point. Since \eqref{LD} defines an irredicuble (every point can be reached) Markov process it must converge to the finite invariant measure $\pi$.
\end{proof}

\begin{theorem}[Exponential convergence]
    Let $\pi$ satisfy the LSI then exponential convergence.
\end{theorem}
\begin{proof}
    By caclulation 
    \begin{aligneq*}
        \partial \KL{\rho_t}{\pi}= - \int_{\Xx} \rho_t \norm{\nabla \log{\frac{\rho_t}{\pi}}}^2 \d x = - I(\rho_t|| \pi). 
    \end{aligneq*}
    By BE $I(\cdot ||\pi)\le \lambda\KL{\cdot }{\pi}$ and we conclude by gronwall.
\end{proof}
\begin{lemma}
   If  $V$ satisfies the Bakry--Emery condition $\nabla^2 V \ge \lambda \rm{Id}$ then $\pi$ verifies the log--Sobolev ineuqality.
\end{lemma}
\begin{proof}
    Since $V$ is strongly convex...
\end{proof}

Convergence may not be exponential id tails decay weakly
\begin{example}
    Consider $V= \bt\log{1+\norm{x}^s}$ where $s>d$ then polynomial convergence. If $V= \log{(1+ \norm{x}^d)\log{\norm{x}}}$ then convergence is logarithmic.
    \exampleSymbol
\end{example}
\section{Proof via WGF}
We define 
\begin{aligneq*}
    \Ff(\rho):= \KL{\rho}{\pi}.
\end{aligneq*}
And we set up the WGF
\begin{aligneq*}
  \partial_t=- \nabla_{W_2} \Ff(\mu_t).
\end{aligneq*}
Since it is a gradient flow we have 
\begin{aligneq*}
  \partial_t \Ff(\mu_t) = - \norm{\nabla_{W_2} \Ff(\mu_t)}_{L^2(\mu_t}^2.
\end{aligneq*}
We have that $\nabla_{W_2} \Ff(\mu)= \nabla \delta \Ff/ \delta \mu= \log{\rho/\pi}$. Plugging this in gives as before that 
\begin{aligneq*}
    \partial_t \Ff(\mu_t)= -I(\mu_t|| \pi).
\end{aligneq*}
We then require for geodesic convexity of $\Ff$ this is equivalent to LSI (I beleive) and shows through gradient flow theory that
\begin{aligneq*}
    \Ff(\mu_t) \le \ex{-2 \lambda t}\Ff(\mu_0).
\end{aligneq*}
And suomehow also that 
\begin{aligneq*}
    W_2(\mu_t,\pi)\le \ex{-2 \lambda t}W_2(\mu_0,\pi)
\end{aligneq*}
\printbibliography
\end{document}
